\begin{claim}{C.82}{Proved}{A null lifted triangulation and its actual skeleton stabilizer}
Retain the actual disk, discrete operator group,
transitive relative subgroup, quotient covering and
metrics of C.81. There is a locally finite
triangulation $\mathcal T$ of $N=O/Q$ whose full
lift $\widetilde{\mathcal T}$ through the original
covering $\pi:O\to N$ has the following property:
for every $\epsilon>0$, only finitely many lifted
closed vertex stars have $d$-diameter at least
$\epsilon$. In particular, this holds for the
lifted closed triangles. Every such star is
mapped injectively by $\pi$.

Write $\Lambda=|\widetilde{\mathcal T}^{(1)}|$
for the actual lifted one-skeleton and set
\[
 \mathscr S_{Q,\Lambda}
   =\{h\in\mathscr C_Q:h|_\Lambda=\id\}.
                                                               \tag{C.82.1}
\]
Choose one closed lifted triangle $D_\sigma$
for each two-simplex $\sigma$ of $\mathcal T$.
Restriction and $Q$-equivariant pasting are
mutually inverse topological group isomorphisms
\[
 \mathscr S_{Q,\Lambda}
    \ \cong\ 
 \prod_{\sigma\in\mathcal T^{(2)}}
       \operatorname{Homeo}(D_\sigma\mathbin{\mathrm{rel}}
                                      \partial D_\sigma).
                                                               \tag{C.82.2}
\]
Here $\mathcal T^{(2)}$ indexes the two-simplices,
the product has the product topology, and the
left side retains the original closed-disk
topology of maps and inverses. Coordinatewise
Alexander deformation gives a continuous
contraction of this entire actual subgroup to
$1$, fixing $1$ throughout.

The conclusion concerns the pointwise skeleton
stabilizer. A deformation of every original
based sphere family in $\mathscr C_Q$ into that
stabilizer, with joint original-boundary control,
is not supplied by this statement or by the
finite-label comparisons of C.81.
\end{claim}

\begin{proof}
\emph{The quotient charts and the chosen subdivision.}
The quotient $N$ is a second-countable surface,
possibly nonorientable. Indeed the original
free proper discrete action gives the actual
covering charts, and the open quotient of the
countable base of $O$ is a countable base of $N$.
Rad\'o's original triangulation theorem,
\cite[\S3, Hilfssatz 2, pp.~111--114]{C_RadoTriangulation},
gives a locally finite triangulation of this
surface. Its proof explicitly establishes
local finiteness before subdividing the finite
polygons into triangles. We retain the
homeomorphism from its geometric realization
onto the actual $N$.

Choose a continuous proper function
$a:N\to[0,\infty)$. To recall its construction,
take compact sets $K_j\subset\operatorname{int}K_{j+1}$
exhausting $N$, and continuous functions
$\varphi_j$ equal to zero on $K_j$ and one
outside $\operatorname{int}K_{j+1}$, with values
in $[0,1]$. The locally finite sum
$a=\sum_j\varphi_j$ is proper: outside $K_{m+1}$
its first $m$ summands equal one. In the compact
case the constant function zero suffices.

For each $y\in N$ choose a relatively compact
open coordinate disk $V_y$ containing $y$ so
small that every sheet above $V_y$ is mapped
isometrically onto $V_y$ for the metrics
$d_Q,\rho$ of C.81, and
\[
 \begin{gathered}
 |a(z)-a(y)|<1\quad(z\in V_y),\\
 \operatorname{diam}_{\rho}(V_y)<\frac1{2+a(y)}.
 \end{gathered}
                                                               \tag{C.82.3}
\]
Such disks are obtained by shrinking the
actual isometric covering neighborhoods of
C.81. No coordinate on the original circle
is changed in this step.

Apply Haver's subdivision lemma with $k=0$
to the initial triangulation, its actual map
to $N$, and the cover $\{V_y\}$
\cite[Lemma 1, p.~458]{C_Haver1975}.
It gives a locally finite subdivision
$\mathcal T$ for which every closed vertex
star is contained in one $V_y$. The notation
$N(v,\mathcal T)$ in that lemma is the union
of the closed simplices meeting the vertex,
so it is precisely the closed star used here.

Lift this actual triangulation through $\pi$.
On each covering chart the lifted simplices
and all their incidences are the pullbacks
under its inverse section. On overlapping
charts these pullbacks agree by covering
uniqueness. They therefore give a triangulation
$\widetilde{\mathcal T}$ of $O$, with the
original $Q$-action by simplicial deck maps.
Every lifted closed vertex star lies in one
sheet over a $V_y$, so $\pi$ restricts to a
homeomorphism on that star. If its quotient
vertex is $v$, then (C.82.3) gives
\[
 \operatorname{diam}_{d_Q}
       \bigl(\overline{\operatorname{st}}(\widetilde v)
                                             \bigr)
       <\frac1{1+a(v)}.
                                                               \tag{C.82.4}
\]
The lift is locally finite, directly from
the covering charts and local finiteness
downstairs. Alternatively, a compact subset
of $O$ projects to a compact set meeting
finitely many quotient stars; properness
allows only finitely many translates of
their compact chosen lifts to meet it.

\emph{Nullity in the original closed disk.}
We use the precise compact-core conclusion
of C.81. For a compact set $C\subset O$,
choose a compact family $T$ of actual
boundary-fixed commuting homeomorphisms
with $C\subset T(x_0)$. Such a family exists
by the compact-covering lemma applied to
the open evaluation map of the closure of
the transitive subgroup. For $z,z'\in C$
choose $s,t\in T$ with $z=sx_0$, $z'=tx_0$.
For every original operator $q\in Q$,
\[
           qz'=t s^{-1}(qz),
       \qquad qC\subset TT^{-1}(qz).
                                                               \tag{C.82.5}
\]
The compact family $TT^{-1}$ consists of
actual maps fixing every point of the
original circle. Hence
\[
 z_j\in C,\quad q_jz_j\longrightarrow b\in\partial\overline D
       \quad\Longrightarrow\quad
            q_jC\longrightarrow\{b\}
             \text{ uniformly in }d.
                                                               \tag{C.82.6}
\]
For completeness, failure of uniformity
would give $h_j\in TT^{-1}$ with
$d(h_jq_jz_j,b)$ bounded away from zero.
A convergent subsequence of $h_j$ in the
actual map topology contradicts $h_jb=b$.
No extension of $q_j$ across the circle
appears in this argument.

Suppose infinitely many distinct lifted
closed stars have diameter at least
$\epsilon>0$. Choose a point $x_j$ in each.
Compactness of $\overline D$ gives a
convergent subsequence. Its limit cannot
lie in $O$, because the lifted stars are
locally finite there. Thus $x_j\to b$ on
the original circle.

Let $v_j$ be their quotient vertices.
An unbounded subsequence of $a(v_j)$
has a further subsequence tending to
infinity. Equation (C.82.4), together
with $d\leq d_Q$, contradicts the lower
diameter bound on that subsequence.
Otherwise the $v_j$ lie in one compact
sublevel of $a$. Local finiteness gives
only finitely many such vertices.
After taking a subsequence, the stars
are therefore $q_jC$ for one chosen
compact lifted star $C$.
Writing $x_j=q_jz_j$ with $z_j\in C$,
(C.82.6) gives the same contradiction.
This proves the asserted finite-large-star
property.

Every closed triangle is contained in
each of its three closed vertex stars.
Only finitely many triangles belong to
any one star. The finite-large-star
property consequently implies the
finite-large-triangle property. It also
gives the exact boundary form that will
be used below: for any lifted triangles
$E_j$ and points $x_j\in E_j$,
\[
 x_j\longrightarrow b\in\partial\overline D
       \quad\Longrightarrow\quad
       \sup_{z\in E_j}d(z,b)\longrightarrow0.
                                                               \tag{C.82.7}
\]
Indeed each fixed finite list of compact
triangles is eventually avoided, and all
others have arbitrarily small diameter.

\emph{The actual product and both inverse maps.}
Choose one lift $D_\sigma$ for each
quotient triangle. A deck transformation
preserving $D_\sigma$ is the identity:
$\pi$ is injective there, so that
transformation fixes its points, and
the original action is free.

For a tuple $(h_\sigma)$ of the factors
in (C.82.2), define $\mathcal P(h_\sigma)$
on $\overline D$ by
\[
 \begin{cases}
   q h_\sigma q^{-1}&\text{on }qD_\sigma,
                           \quad q\in Q,\\
   \id&\text{on }\Lambda\cup\partial\overline D.
 \end{cases}
                                                               \tag{C.82.8}
\]
The first prescription is unambiguous
on every open face. On common edges
and vertices all prescriptions are
the identity. Local finiteness gives
continuity on $O$ by finite pasting.
Equation (C.82.7) gives continuity at
the original circle: the image of a
point in a triangle stays in that same
triangle. Replacing every $h_\sigma$
by its inverse gives a continuous map
by exactly the same proof, and the two
composites are the identity face by
face and on the skeleton and circle.
Thus (C.82.8) is an actual closed-disk
homeomorphism. Reindexing its formula
by left multiplication in $Q$ proves
commutation with every original operator.

Conversely, an element $h$ of (C.82.1)
preserves each closed lifted triangle.
Its Jordan boundary is fixed pointwise.
Of the two complementary components
in $O$, only its interior has compact
closure in $O$. An interior homeomorphism
preserves this property, so it cannot
exchange the two components. Restriction
therefore defines
\[
 \mathcal R(h)=(h|_{D_\sigma})_\sigma.
                                                               \tag{C.82.9}
\]
The commutator equations imply the
formula (C.82.8) on every other lift.
It follows that $\mathcal P\mathcal R$
and $\mathcal R\mathcal P$ are the
identity. Both maps respect the group
operations. This proves the asserted
group isomorphism with its exact maps.

Restriction $\mathcal R$ is continuous
for maps and inverses, since each
$D_\sigma$ is compact and invariant.
For continuity of $\mathcal P$, fix
$\epsilon>0$. Only finitely many actual
lifted triangles have diameter at least
$\epsilon/4$. On every other triangle,
two pasted maps, and separately their
inverses, take each point into the same
set of diameter less than $\epsilon/4$.
On the finitely many large triangles,
uniform convergence of the corresponding
finitely many factors gives uniform
convergence after their specified
conjugations by $q$. This uses continuity
of that fixed $q$ on a compact triangle,
and requires no boundary extension of
$q$. The finitely many factor neighborhoods
therefore control both uniform distances
on all of $\overline D$. This is exactly
continuity from the product topology
to the original map-and-inverse topology.
It proves (C.82.2) as a topological
group isomorphism.

\emph{A contraction in the actual topology.}
Choose a homeomorphism from the standard
closed unit disk to each $D_\sigma$.
In that coordinate, for a boundary-fixed
disk homeomorphism $h$ put
\[
 A_r(h)(z)=
 \begin{cases}
      r h(z/r),&|z|\leq r,\quad r>0,\\
      z,&|z|\geq r,
 \end{cases}
 \qquad A_0(h)=\id.
                                                               \tag{C.82.10}
\]
The two formulas agree at $|z|=r$,
and $A_r(h)^{-1}=A_r(h^{-1})$.
For positive $r$, evaluation is jointly
continuous, including across that seam;
compactness of the standard disk gives
the required compact-open continuity.
At $r=0$ the map and inverse displacements
are each at most $2r$, uniformly in $h$.
This proves joint continuity for the
full map-and-inverse group topology.
Also $A_1(h)=h$ and $A_r(\id)=\id$.

Take the product of these deformations
and transport it through the mutually
inverse maps $\mathcal P,\mathcal R$.
A map into a product is continuous
exactly when its coordinate maps are
continuous. Thus this gives a continuous
map $[0,1]\times\mathscr S_{Q,\Lambda}
\to\mathscr S_{Q,\Lambda}$ from the
constant identity map to the identity
map, fixing the group identity throughout.
All original operator labels, the
pointwise skeleton condition, and the
original circle condition remain exact.

The construction proves an actual
infinite product comparison because
the nullity estimate supplies uniform
control outside finitely many triangles.
It does not assert an infinite-configuration
lifting theorem for restriction from
$\mathscr C_Q$. In particular, C.81's
finite-orbit evaluation bundles have
not been passed to an inverse limit.
The deformation of a general compact
parameter family to this skeleton
stabilizer, with controlled vertex and
edge motions, remains the next step
in the discrete-operator part of C.12.
\end{proof}
